\begin{lemma}[Common constants and uniform order bounds]
Put $Q=\Delta^2+1$ and let $\mathcal K$ be an induction application
certificate with parameters
$(\Delta,Q,d,\theta,p,\varepsilon,\rho,\zeta,\omega,\widehat\zeta)$
in the sense of \noderef{InductionApplicationCertificate}.
Assume that for $s<Q$,
\[
 0<\widehat\zeta_s<\zeta_s<\infty,\qquad
 \omega_{s+1}<\varepsilon_{s+1}\leq1,
\]
and every $\mathcal K$-matching constant at $s<Q$, $x\leq\Delta$ is
finite. Then there are $0<a\leq1\leq M$, $C\geq2M$, and
$n_0\in\mathbb N$ such that, for every $s<Q$ and $x,y\leq\Delta$,
\[
 a\leq\mathcal K.\mathrm{fourGraphsConstant}(s,x,y),\qquad
 \mathcal K.\mathrm{matchingConstant}(s,x)\leq M,
\]
and each witness-thinning, host-thinning, matching-spread,
matching-existence, and Four Graphs threshold at these indices is at
most $n_0$. For every $n\geq n_0$ and $s<Q$,
\[
 C/n\leq1,\qquad \Delta(1-a)^n\leq\tfrac12,
\]
\[
 n^{-1}\leq\widehat\zeta_s,\quad
 n^{-1}\leq\zeta_s-\widehat\zeta_s,\quad
 \widehat\zeta_s+2/n\leq\zeta_s,\quad
 \widehat\zeta_s+1/(\zeta_s n)\leq\zeta_s,
\]
\[
 \omega_{s+1}+2/n\leq\varepsilon_{s+1},\qquad
 \omega_{s+1}+1/(\varepsilon_{s+1}n)\leq\varepsilon_{s+1}.
\]
\end{lemma}
\begin{proof}
Let $S=\{0,\ldots,Q-1\}$ and $X=\{0,\ldots,\Delta\}$.
Both are finite and nonempty, also when $\Delta=0$.
By the constant fields of \noderef{InductionApplicationCertificate},
every $a_{s,x,y}=\mathcal K.\mathrm{fourGraphsConstant}(s,x,y)$ is
positive and at most one, and every
$M_{s,x}=\mathcal K.\mathrm{matchingConstant}(s,x)$ is at least one.
The latter are finite by hypothesis. Set
\[
 a=\min\bigl(\{1\}\cup\{a_{s,x,y}:s\in S,\ x,y\in X\}\bigr),
 \quad
 M=\max\bigl(\{1\}\cup\{M_{s,x}:s\in S,\ x\in X\}\bigr),
 \quad C=2M.
\]
The finite minimum is positive and at most one, and the finite maximum
is finite and at least one. Thus $a,M,C$ are finite, and they satisfy
all required constant comparisons, including those with $y=x$.
Let $N_{\rm cert}$ be the maximum of zero and all the finitely many
certificate thresholds in the asserted index ranges. Each such
threshold is a natural number, so this maximum exists.

For each $s\in S$ set
\[
 r_s=\zeta_s-\widehat\zeta_s,\qquad
 q_s=\varepsilon_{s+1}-\omega_{s+1}.
\]
Since $\widehat\zeta_s<\zeta_s<\infty$, $r_s$ is positive finite;
also $\zeta_s$ and $\widehat\zeta_s$ are positive finite.
The other strict inequality gives
$0<\varepsilon_{s+1}\leq1$ and
$0\leq\omega_{s+1}<\varepsilon_{s+1}<\infty$, hence $q_s$ is positive
finite. All subtractions here therefore agree with ordinary real
subtraction, and the products $\zeta_s r_s$ and
$\varepsilon_{s+1}q_s$ are positive finite. In particular all the real
numbers in the following finite set are finite:
\[
 \{1,C\}\cup\bigcup_{s\in S}
 \left\{\frac1{\widehat\zeta_s},\frac1{r_s},
 \frac2{r_s},\frac1{\zeta_s r_s},
 \frac2{q_s},\frac1{\varepsilon_{s+1}q_s}\right\}.
\]
Use the Archimedean property to choose a natural number $N_{\rm pad}$
at least every member of this set. We identify finite extended real
numbers with their real values only after the finiteness checks above.
For $n\geq N_{\rm pad}$, $n\geq1$ and $n\geq C$, so $C/n\leq1$.
For each $s\in S$, multiplication or division by the indicated positive
finite denominators gives
\[
 1/n\leq\widehat\zeta_s,\quad1/n\leq r_s,
 \quad2/n\leq r_s,\quad1/(\zeta_s n)\leq r_s,
 \quad2/n\leq q_s,\quad
 1/(\varepsilon_{s+1}n)\leq q_s.
\]
Adding $\widehat\zeta_s$ to the third and fourth inequalities gives
the two inner padding bounds; adding $\omega_{s+1}$ to the last two
gives the two output padding bounds. The first inequality is the
explicit inner reserve, independent of the second inequality about the
gap $r_s$; neither has been inferred from the other.

It remains to bound the geometric failure term. If $\Delta=0$, it is
zero for every $n$. If $a=1$, it is zero for every $n\geq1$.
Otherwise $0<a<1$. Write $b=1-a\in(0,1)$ and
$u=a/(1-a)>0$. An elementary induction gives
$(1+u)^n\geq1+nu$: the induction step multiplies by $1+u$ and
uses the nonnegative additional term $nu^2$.
Since $1+u=1/b$ and $u\geq a$, inversion gives
\[
 b^n\leq\frac1{1+nu}\leq\frac1{1+na}.
\]
Choose an integer $N_{\rm decay}\geq2\Delta/a$ by the Archimedean
property, which applies because $a>0$ is finite. For every
$n\geq N_{\rm decay}$ one has
$2\Delta\leq na\leq1+na$, so
\[
 \Delta(1-a)^n\leq\frac{\Delta}{1+na}\leq\frac12.
\]
In the two immediate cases take $N_{\rm decay}=1$.

Finally set $n_0=\max(N_{\rm cert},N_{\rm pad},N_{\rm decay})$.
Every certificate threshold is at most $n_0$. For every $n\geq n_0$
the preceding arguments establish each of the claimed inequalities.
All choices are functions only of the finite parameter and certificate
data and $\Delta$; no reduced-graph order or instance occurs.
\end{proof}
